\begin{abstract}
Execution-free verification promises to replace costly test execution in Large Language Model (LLM)-based program repair, yet practitioners must choose among three fundamentally different formulations with no controlled comparison to guide the decision: per-test classification (\cls{}), execution output generation (\cwm{}), and trajectory-level scoring (\traj{}). We present the first systematic comparison, controlling for model capacity, training data, and evaluation protocol across two code LLM families (Qwen3, DeepSeek-Coder), two complexity levels, and multiple random seeds under ${\leq}$8B Low-Rank Adaptation (LoRA) fine-tuning. Our central finding is that \emph{code complexity moderates the classification--generation tradeoff}. When outputs are short (function-level verification), classification and generation achieve statistically equivalent accuracy; trajectory-level parity is model-dependent. When outputs grow longer and more structured (repository-level verification), this equivalence breaks down: generation-based verification suffers format collapse (failure to produce parseable outputs), falling well below classification. At 14B, function-level gaps narrow to approximately one percentage point, but whether the repo-level divergence is capacity-contingent remains open. A failure mode taxonomy and error complementarity analysis provide actionable guidance for formulation selection and motivate complexity-aware routing as a key open direction\footnote{Code and data are available at \url{https://anonymous.4open.science/r/exec-sim-repair-verifier-102B/}.}.
\end{abstract}
